\documentclass[12pt]{article}
\usepackage[margin=1in]{geometry}
\usepackage{amsmath,amssymb,amsthm}
\usepackage{booktabs}
\usepackage[colorlinks=true,linkcolor=blue,citecolor=blue,urlcolor=blue]{hyperref}

\newcommand{\seqnum}[1]{\href{https://oeis.org/#1}{\underline{#1}}}
\newcommand{\fl}[1]{\lfloor #1 \rfloor}
\newcommand{\B}{\mathcal{B}}
\newcommand{\Z}{\mathbb{Z}}
\newcommand{\Q}{\mathbb{Q}}

\theoremstyle{plain}
\newtheorem{theorem}{Theorem}
\newtheorem{lemma}[theorem]{Lemma}
\newtheorem{proposition}[theorem]{Proposition}
\newtheorem{corollary}[theorem]{Corollary}
\theoremstyle{remark}
\newtheorem{remark}[theorem]{Remark}

\title{Gaps in the Intersection of Two Beatty Sequences:\\ Three Conjectures of Kimberling}
\author{Camden Link}
\date{Draft of November 21, 2026}

\begin{document}
\maketitle

\begin{abstract}
The Beatty sequence of an irrational number greater than one consists of the integer parts of
its positive multiples. Fraenkel and Holzman proved that the integers common to two Beatty
sequences have only finitely many distinct gaps between consecutive elements, by an argument
that does not compute them. We determine the set of gaps, and the frequency of each gap, for
four pairs of Beatty sequences recorded in the On-Line Encyclopedia of Integer Sequences by
Kimberling. For the pair formed by the square root of three and one plus the square root of
three, this settles the three conjectures in the corresponding entry: the gap set is not the
conjectured one, the growth rate is as conjectured, and the sequence returns to its linear trend
infinitely often while deviating from it without bound in both directions. The last statement
follows from a description of the deviation as a walk with unit steps driven by the gap
sequence. The gap sets are formally verified in the Lean proof assistant, and so is the last
statement, conditionally on one classical fact about Sturmian words.
\end{abstract}

\section{Introduction}

For an irrational number $\alpha>1$, the \emph{Beatty sequence} of $\alpha$ is the increasing
sequence $\fl{\alpha},\fl{2\alpha},\fl{3\alpha},\dots$; we write
$\B(\alpha)=\{\fl{n\alpha} : n\ge 1\}$ for its set of values. Given irrational
$\alpha,\beta>1$ we study the intersection $\B(\alpha)\cap\B(\beta)$. When it is infinite we
list its elements as $a(1)<a(2)<\cdots$ and call the differences $a(n+1)-a(n)$ its
\emph{gaps}.

Skolem \cite{skolem} showed that $\B(\alpha)\cap\B(\beta)$ is infinite when $1,1/\alpha,1/\beta$
are linearly independent over $\Q$, and in some dependent cases. Bang \cite{bang} determined
its density; in the form quoted by Begunts and Goryashin \cite[Theorem~1]{bg25}, the density is
$1/(\alpha\beta)$ in the independent case, and if $u/\alpha+v/\beta=w$ for integers $u,v,w$
with greatest common divisor $1$, it is
\begin{equation}\label{eq:bang}
\delta=\frac{1}{\alpha\beta}-\frac{1}{uv}\Bigl\{\frac{u}{\alpha}\Bigr\}\Bigl\{\frac{v}{\beta}\Bigr\},
\end{equation}
where $\{x\}$ denotes the fractional part. Begunts and Goryashin \cite{bg22,bg25} bound the
error term in both cases. The gaps were first studied by Fraenkel and Holzman \cite{fh}, who
reduced the question to return times of a translation of the two-dimensional torus to a
rectangle and proved that the number of distinct gaps is always finite
\cite[Corollary~5.2]{fh}. Their proof is by compactness and gives neither the gaps nor a bound.
For rational moduli they obtained sharp bounds; see also Simpson \cite{simpson}.

In 2021--22 Kimberling entered several such intersections in the OEIS \cite{oeis}, with
conjectures. The entry \seqnum{A352676}, the intersection for $\sqrt3$ and $1+\sqrt3$, states:
\begin{quote}
(1) $a(n+1)-a(n)$ is in $\{2,3,4,5,6\}$ for every $n$, and each of these differences occurs
infinitely many times.\\
(2) $\lim_{n\to\infty} a(n)/n = (3/2)(1+\sqrt3)$.\\
(3) Let $d(n) = a(n) - \seqnum{A352673}(n)$; then $d(n) = 0$ for infinitely many $n$, but
$\{d(n)\}$ is unbounded below and above.
\end{quote}
The sequence \seqnum{A352673} is the decimal expansion of a constant, so (3) cannot be meant
literally. We read it with \seqnum{A352674}, the Beatty sequence of $(3/2)(1+\sqrt3)$, in place
of \seqnum{A352673}; this matches the parallel conjecture in \seqnum{A346308}, where
$d(n)=a(n)-\fl{n\sqrt6}$. The entry \seqnum{A351415}, the intersection for the golden ratio and
$\sqrt5$, conjectures that every gap lies in $\{2,3,4,5,6\}$ and that each of these occurs
infinitely often. The entries \seqnum{A346308} (for $\sqrt2$ and $\sqrt3$) and
\seqnum{A347793} (for $2^{1/3}$ and $2^{2/3}$) carry conjectures of type (3) only.

\subsection*{Results}

By the \emph{frequency} of a gap $g$ we mean the limit, as $x\to\infty$, of the proportion of
$n\le x$ with $a(n+1)-a(n)=g$.

\begin{theorem}\label{thm:gaps}
For each of the following pairs, the set of gaps of $\B(\alpha)\cap\B(\beta)$ and the
frequencies of the gaps are as shown. In particular every listed gap occurs infinitely often.
\begin{center}
\renewcommand{\arraystretch}{1.25}
\begin{tabular}{llll}
\toprule
$\alpha,\ \beta$ & OEIS & gaps & frequencies, in order\\
\midrule
$\sqrt3,\ 1+\sqrt3$ & \seqnum{A352676} & $2,3,5,6$ &
 $\frac{2-\sqrt3}{2},\ \frac{\sqrt3-1}{2},\ 2-\sqrt3,\ \sqrt3-\frac32$\\
$\frac{1+\sqrt5}{2},\ \sqrt5$ & \seqnum{A351415} & $2,3,4,5,6$ &
 $\frac{7+\sqrt5}{22},\ \frac{4-\sqrt5}{11},\ \frac{17-7\sqrt5}{11},\ \frac{5\sqrt5-9}{11},\ \frac{5\sqrt5-9}{22}$\\
$\sqrt2,\ \sqrt3$ & \seqnum{A346308} & $1,2,3,4$ &
 $(2-\sqrt2)(2-\sqrt3),\ (2-\sqrt2)(\sqrt3-1),$\\
 & & & $2(\sqrt2-1)(2-\sqrt3),\ (\sqrt2-1)(2\sqrt3-3)$\\
$\lambda,\ \lambda^2$ with $\lambda=2^{1/3}$ & \seqnum{A347793} & $1,2,3,4$ &
 $2(\lambda^2-1)(\lambda-1),\ (\lambda^2-1)(2-\lambda),$\\
 & & & $2(2-\lambda^2)(\lambda-1),\ \tfrac12(2-\lambda^2)(4-3\lambda)$\\
\bottomrule
\end{tabular}
\end{center}
\end{theorem}

\begin{theorem}\label{thm:a352676}
Let $a(n)$ be the $n$-th element of $\B(\sqrt3)\cap\B(1+\sqrt3)$.
\begin{itemize}
\item[(i)] Conjecture (1) of \seqnum{A352676} is false as stated: the gap $4$ never occurs. The
gaps are exactly $2,3,5,6$, and each occurs infinitely often.
\item[(ii)] Conjecture (2) holds: $a(n)/n\to(3/2)(1+\sqrt3)$.
\item[(iii)] Conjecture (3) holds in the reading above: with
$d(n)=a(n)-\fl{(3/2)(1+\sqrt3)\,n}$, we have $d(n)=0$ for infinitely many $n$, and $d$ is
unbounded above and below.
\end{itemize}
\end{theorem}

Part (ii) is a consequence of Bang's formula \eqref{eq:bang} with $(u,v,w)=(3,-2,1)$; a
recipe that yields it was also given by Kalmynin \cite{kalmynin} in answer to a question of
Kimberling. We include a proof because it costs one line once the gaps are understood. Part
(iii) rests on the following description, which we find the most interesting result of the
paper. Call a gap \emph{short} if it is $2$ or $3$ and \emph{long} if it is $5$ or $6$.

\begin{theorem}\label{thm:walk}
With $a(n)$ and $d(n)$ as in Theorem~\ref{thm:a352676}, let $G(n)$ be $-1$ plus the number of
short gaps minus the number of long gaps among the first $n-1$ gaps
$a(2)-a(1),\dots,a(n)-a(n-1)$. Then:
\begin{itemize}
\item[(a)] $d(n)=-1-\bigl\lfloor (G(n)-\varphi_n)/(\sqrt3-1)\bigr\rfloor$, where
$\varphi_n=\{(a(n)+1)\sqrt3\}$ satisfies $0<\varphi_n<\sqrt3-1$. In particular $d(n)=0$ if and
only if $G(n)=0$, and $d(n)$ and $G(n)$ have opposite signs otherwise.
\item[(b)] Let $q_0=q_1=1$ and $q_k=2q_{k-1}+q_{k-2}$ for even $k\ge2$,
$q_k=q_{k-1}+q_{k-2}$ for odd $k\ge3$. For $j\ge1$, if $a(n)\le q_{4j+1}<a(n+1)$ then
$2j-1\le G(n)\le 2j+2$, and if $a(n)\le q_{4j+3}<a(n+1)$ then $-2j-3\le G(n)\le -2j$.
\end{itemize}
\end{theorem}

Since $G$ changes by exactly one at each step, (b) forces it to take the value $0$ infinitely
often, and (a) then gives Theorem~\ref{thm:a352676}(iii). The numbers $q_k$ are the denominators
$1,1,3,4,11,15,41,56,\dots$ of the convergents of $\sqrt3-1$; they grow geometrically, so
$|d(n)|$ exceeds a constant multiple of $\log n$ infinitely often on each side.

\begin{remark}\label{rem:four}
That $4$ is not a gap in Theorem~\ref{thm:a352676}(i) has a two-line proof: $4$ is not even a
difference of two elements of $\B(1+\sqrt3)$. Indeed, with $\beta=1+\sqrt3$, consecutive
elements differ by $2$ or $3$, and $\fl{(k+2)\beta}-\fl{k\beta}>2\beta-1>4$. The substance of
(i) is that $2,3,5,6$ all occur and nothing else does.
\end{remark}

\subsection*{Formal verification}

The gap sets in Theorem~\ref{thm:gaps} (that the gaps are exactly those listed, and that each
occurs at least once) are formally verified in Lean~4 \cite{lean4} with Mathlib \cite{mathlib}.
So is Theorem~\ref{thm:a352676}(iii), under one explicitly stated hypothesis: the classical fact
about standard words quoted in Section~\ref{sec:drift}. The frequencies are proved here in the
usual way and are not formalized.
Section~\ref{sec:lean} says precisely what is and is not machine-checked. The Lean sources and
the programs that produced the case analyses are available at
\url{https://github.com/clink-bw/math-proofs}.

\subsection*{Organization}

Section~\ref{sec:prelim} reduces membership in a Beatty sequence to a statement about one real
offset. Sections~\ref{sec:sqrt3}--\ref{sec:indep} determine the gap sets.
Section~\ref{sec:drift} proves Theorem~\ref{thm:walk}. Section~\ref{sec:lean} describes the
formalization, Section~\ref{sec:open} lists open questions, and Section~\ref{sec:ai} states how
AI tools were used.

\section{Preliminaries}\label{sec:prelim}

Throughout, $\alpha>1$ is irrational and $m\ge1$ is an integer.

\begin{lemma}\label{lem:member}
We have $m\in\B(\alpha)$ if and only if $\fl{(m+1)/\alpha}>\fl{m/\alpha}$.
\end{lemma}

\begin{proof}
$m=\fl{k\alpha}$ means $m<k\alpha<m+1$ (equality is impossible), that is,
$m/\alpha<k<(m+1)/\alpha$. Such an integer $k$ exists if and only if
$\fl{(m+1)/\alpha}>\fl{m/\alpha}$, because $(m+1)/\alpha$ is not an integer; and then $k\ge1$.
\end{proof}

\begin{lemma}\label{lem:hit}
Let $\sigma>0$ be such that $d=\alpha\sigma$ is a positive integer. Then $m\in\B(\alpha)$ if
and only if there is an integer $k$ with
\begin{equation}\label{eq:hit}
m\sigma\le dk<m\sigma+\sigma .
\end{equation}
The integer $k$ is then unique, and we call $dk-m\sigma\in[0,\sigma)$ the \emph{offset} of $m$.
\end{lemma}

\begin{proof}
Multiply $m\le k\alpha<m+1$ by $\sigma$. Uniqueness holds because $\sigma<d$.
\end{proof}

\begin{lemma}\label{lem:shift}
In the situation of Lemma~\ref{lem:hit}, let $m\in\B(\alpha)$ have offset $a$, and let $j\ge1$.
Then $m+j\in\B(\alpha)$ if and only if
\[
a\in\bigcup_{i\in\Z}\bigl[\,j\sigma-di,\ (j+1)\sigma-di\,\bigr).
\]
\end{lemma}

\begin{proof}
Write $a=dk-m\sigma$. By \eqref{eq:hit}, $m+j\in\B(\alpha)$ if and only if
$(m+j)\sigma\le d(k+i)<(m+j)\sigma+\sigma$ for some integer $i$, that is,
$j\sigma\le a+di<(j+1)\sigma$.
\end{proof}

We use Lemma~\ref{lem:hit} with the following choices.
\begin{center}
\begin{tabular}{llll}
\toprule
$\alpha$ & $\sigma$ & $d$ & reason\\
\midrule
$\sqrt D$ for $D=2,3,5$ & $\sqrt D$ & $D$ & \\
$1+\sqrt3$ & $\sqrt3-1$ & $2$ & $(1+\sqrt3)(\sqrt3-1)=2$\\
$(1+\sqrt5)/2$ & $\sqrt5-1$ & $2$ & $(1+\sqrt5)(\sqrt5-1)=4$\\
$\lambda=2^{1/3}$ & $\lambda^2$ & $2$ & \\
$\lambda^2$ & $\lambda$ & $2$ & \\
\bottomrule
\end{tabular}
\end{center}

For frequencies we need Weyl's theorem \cite{weyl,kn}: if $\omega$ is irrational then the
fractional parts $\{m\omega\}$ are equidistributed in $[0,1)$, and if $1,\omega_1,\omega_2$ are
linearly independent over $\Q$ then the points $(\{m\omega_1\},\{m\omega_2\})$ are
equidistributed in the unit square.

\section{The pair \texorpdfstring{$\sqrt3$}{sqrt 3} and \texorpdfstring{$1+\sqrt3$}{1 + sqrt 3}}\label{sec:sqrt3}

Put $\theta=\sqrt3$ and $S=\B(\theta)\cap\B(1+\theta)$. For $m\in S$, Lemma~\ref{lem:hit} gives
integers $k,l$ and offsets
\[
a=3k-m\theta\in[0,\theta),\qquad b=2l-m(\theta-1)\in[0,\theta-1).
\]
Their difference $c=b-a=2l+m-3k$ is an integer with $-\theta<c<\theta-1$, so $c\in\{0,-1\}$.
By Lemma~\ref{lem:shift}, for $j\ge1$,
\[
m+j\in\B(\theta)\iff a\in P_j,\qquad m+j\in\B(1+\theta)\iff b\in Q_j,
\]
where $P_j=\bigcup_i[j\theta-3i,(j+1)\theta-3i)$ and
$Q_j=\bigcup_i[j(\theta-1)-2i,(j+1)(\theta-1)-2i)$. Intersecting with the ranges of $a$ and $b$
gives Table~\ref{tab:sqrt3PQ}; for instance $P_3$ meets $[0,\theta)$ only in its piece $i=2$,
which is $[3\theta-6,4\theta-6)$.

\begin{table}[ht]
\centering
\begin{tabular}{ccc}
\toprule
$j$ & $P_j\cap[0,\theta)$ & $Q_j\cap[0,\theta-1)$\\
\midrule
1 & $[0,\,2\theta-3)$ & $\emptyset$\\
2 & $[2\theta-3,\,\theta)$ & $[0,\,3\theta-5)$\\
3 & $[0,\,4\theta-6)$ & $[3\theta-5,\,\theta-1)$\\
4 & $[4\theta-6,\,\theta)$ & $\emptyset$\\
5 & $[0,\,6\theta-9)$ & $[0,\,6\theta-10)$\\
6 & $[6\theta-9,\,\theta)$ & $[6\theta-10,\,\theta-1)$\\
\bottomrule
\end{tabular}
\caption{The sets $P_j$ and $Q_j$ for $\theta=\sqrt3$. Numerically $2\theta-3\approx0.464$,
$4\theta-6\approx0.928$, $6\theta-9\approx1.392$, $3\theta-5\approx0.196$,
$6\theta-10\approx0.392$.}\label{tab:sqrt3PQ}
\end{table}

\begin{proposition}\label{prop:sqrt3cells}
Let $m\in S$ have offsets $a$, $b=a+c$. The next element of $S$ after $m$ is $m+g$, where $g$ is
given by Table~\ref{tab:sqrt3cells}.
\end{proposition}

\begin{table}[ht]
\centering
\begin{tabular}{cllccc}
\toprule
$c$ & range of $a$ & range of $b$ & $g$ & length & $\fl{\theta-a},\ \fl{(g+1)\theta-a}$\\
\midrule
$0$ & $[0,\,3\theta-5)$ & same & 5 & $3\theta-5$ & $1,\ 10$\\
$0$ & $[3\theta-5,\,2\theta-3)$ & same & 3 & $2-\theta$ & $1,\ 6$\\
$0$ & $[2\theta-3,\,\theta-1)$ & same & 3 & $2-\theta$ & $1,\ 6$\\
$-1$ & $[1,\,3\theta-4)$ & $[0,\,3\theta-5)$ & 2 & $3\theta-5$ & $0,\ 4$\\
$-1$ & $[3\theta-4,\,6\theta-9)$ & $[3\theta-5,\,6\theta-10)$ & 5 & $3\theta-5$ & $0,\ 9$\\
$-1$ & $[6\theta-9,\,\theta)$ & $[6\theta-10,\,\theta-1)$ & 6 & $9-5\theta$ & $0,\ 10$\\
\bottomrule
\end{tabular}
\caption{The six cells for $\sqrt3$ and $1+\sqrt3$. The last column is used in
Section~\ref{sec:drift}.}\label{tab:sqrt3cells}
\end{table}

\begin{proof}
The next element is $m+g$ for the least $j=g$ with $a\in P_j$ and $b\in Q_j$. Since
$Q_1=Q_4=\emptyset$, only $j\in\{2,3,5,6\}$ can occur. We read off Table~\ref{tab:sqrt3PQ}.

\emph{Case $c=0$}, so $b=a\in[0,\theta-1)$. If $a<3\theta-5$ then $a\notin P_2$ and $b\notin Q_3$,
while $a\in P_5$ and $b\in Q_5$: $g=5$. If $3\theta-5\le a<2\theta-3$ then $a\notin P_2$,
$a\in P_3$, $b\in Q_3$: $g=3$. If $a\ge2\theta-3$ then $b\notin Q_2$, $a\in P_3$, $b\in Q_3$:
$g=3$.

\emph{Case $c=-1$}, so $a\in[1,\theta)$ and $b=a-1$. If $a<3\theta-4$ then $a\in P_2$ and
$b\in Q_2$: $g=2$. If $3\theta-4\le a<6\theta-9$ then $b\notin Q_2$, $a\notin P_3$, and
$a\in P_5$, $b\in Q_5$: $g=5$. If $a\ge6\theta-9$ then $b\notin Q_2$, $a\notin P_3$,
$a\notin P_5$, and $a\in P_6$, $b\in Q_6$: $g=6$.
\end{proof}

\begin{proof}[Proof of Theorem~\ref{thm:gaps} for \seqnum{A352676} and of
Theorem~\ref{thm:a352676}(i),(ii)]
Let $y_m\in[0,6)$ be the residue of $-m(3+\theta)$ modulo $6$. Then $a\equiv y_m\pmod 3$ and
$b\equiv y_m\pmod 2$, because $-m(3+\theta)\equiv-m\theta\pmod3$ and
$-m(3+\theta)\equiv m-m\theta\pmod 2$. Hence $m\in S$ if and only if $y_m\bmod 3<\theta$ and
$y_m\bmod 2<\theta-1$, which holds exactly for $y_m\in[0,\theta-1)\cup[4,3+\theta)$; on the
first interval $a=y_m$ and $c=0$, on the second $a=y_m-3$ and $c=-1$. So each cell of
Table~\ref{tab:sqrt3cells} is an interval of values of $y_m$ of the stated length. As
$(3+\theta)/6$ is irrational, $(y_m)$ is equidistributed in $[0,6)$. Therefore every cell is
visited infinitely often, the density of $S$ is $2(\theta-1)/6=(\theta-1)/3$, and the
frequency of a gap is the total length of its cells divided by $2(\theta-1)$. The lengths are
$3\theta-5$, $4-2\theta$, $6\theta-10$, $9-5\theta$ for $g=2,3,5,6$, which gives the
frequencies in Theorem~\ref{thm:gaps}. Finally $n/a(n)\to(\theta-1)/3$, that is,
$a(n)/n\to3/(\theta-1)=(3/2)(1+\theta)$.
\end{proof}

Note that the short gaps have total frequency $\frac12$, as do the long gaps.

\section{The golden ratio and \texorpdfstring{$\sqrt5$}{sqrt 5}}\label{sec:golden}

Put $\eta=\sqrt5$ and $S=\B(\eta)\cap\B((1+\eta)/2)$. For $m\in S$ the offsets are
\[
a=5k-m\eta\in[0,\eta),\qquad b=2l-m(\eta-1)\in[0,\eta-1),
\]
and $c=b-a=2l+m-5k$ is an integer in $(-\eta,\eta-1)$, so $c\in\{1,0,-1,-2\}$. Now
$P_j=\bigcup_i[j\eta-5i,(j+1)\eta-5i)$ and $Q_j=\bigcup_i[j(\eta-1)-2i,(j+1)(\eta-1)-2i)$,
which gives Table~\ref{tab:goldenPQ}.

\begin{table}[ht]
\centering
\begin{tabular}{ccc}
\toprule
$j$ & $P_j\cap[0,\eta)$ & $Q_j\cap[0,\eta-1)$\\
\midrule
1 & $\emptyset$ & $[0,\,2\eta-4)$\\
2 & $[0,\,3\eta-5)$ & $[2\eta-4,\,\eta-1)$\\
3 & $[3\eta-5,\,\eta)$ & $[0,\,4\eta-8)$\\
4 & $[0,\,5\eta-10)$ & $[0,\,5\eta-11)\cup[4\eta-8,\,\eta-1)$\\
5 & $[5\eta-10,\,\eta)$ & $[5\eta-11,\,\eta-1)$\\
6 & $[0,\,7\eta-15)$ & $[0,\,7\eta-15)$\\
\bottomrule
\end{tabular}
\caption{The sets $P_j$ and $Q_j$ for $\eta=\sqrt5$. Numerically $3\eta-5\approx1.708$,
$5\eta-10\approx1.180$, $7\eta-15\approx0.652$, $2\eta-4\approx0.472$, $4\eta-8\approx0.944$,
$5\eta-11\approx0.180$.}\label{tab:goldenPQ}
\end{table}

\begin{proposition}\label{prop:goldencells}
Let $m\in S$ have offsets $a$, $b=a+c$. The next element of $S$ after $m$ is $m+g$, where $g$ is
given by Table~\ref{tab:goldencells}.
\end{proposition}

\begin{table}[ht]
\centering
\begin{tabular}{clcc@{\hspace{3em}}clcc}
\toprule
$c$ & range of $a$ & $g$ & length & $c$ & range of $a$ & $g$ & length\\
\midrule
$1$ & $[0,\,\eta-2)$ & 2 & $\eta-2$ & $-1$ & $[5\eta-10,\,2\eta-3)$ & 5 & $7-3\eta$\\
$0$ & $[0,\,5\eta-11)$ & 4 & $5\eta-11$ & $-1$ & $[2\eta-3,\,3\eta-5)$ & 2 & $\eta-2$\\
$0$ & $[5\eta-11,\,2\eta-4)$ & 6 & $7-3\eta$ & $-1$ & $[3\eta-5,\,4\eta-7)$ & 3 & $\eta-2$\\
$0$ & $[2\eta-4,\,\eta-1)$ & 2 & $3-\eta$ & $-1$ & $[4\eta-7,\,\eta)$ & 5 & $7-3\eta$\\
$-1$ & $[1,\,5\eta-10)$ & 4 & $5\eta-11$ & $-2$ & $[2,\,\eta)$ & 3 & $\eta-2$\\
\bottomrule
\end{tabular}
\caption{The ten cells for $\sqrt5$ and the golden ratio; in each, $b=a+c$.}\label{tab:goldencells}
\end{table}

\begin{proof}
As before, $g$ is the least $j$ with $a\in P_j$ and $b\in Q_j$, and $j=1$ is excluded because
$P_1=\emptyset$. Each row is checked against Table~\ref{tab:goldenPQ}. We give two rows; the
others are similar. For $c=0$ and $5\eta-11\le a<2\eta-4$: $b=a\notin Q_2$, $a\notin P_3$,
$b\notin Q_4$, $a\notin P_5$, and $a,b<7\eta-15$, so $g=6$. For $c=-1$ and
$4\eta-7\le a<\eta$, so that $b=a-1\in[4\eta-8,\eta-1)$: $a\notin P_2$, $b\notin Q_3$,
$a\notin P_4$, and $a\in P_5$, $b\in Q_5$, so $g=5$.
\end{proof}

\begin{proof}[Proof of Theorem~\ref{thm:gaps} for \seqnum{A351415}]
Let $y_m\in[0,10)$ be the residue of $-m(5+\eta)$ modulo $10$; then $a\equiv y_m\pmod5$ and
$b\equiv y_m\pmod2$, and $(y_m)$ is equidistributed in $[0,10)$. As in
Section~\ref{sec:sqrt3}, the cells are intervals of values of $y_m$ of the stated lengths. The
total lengths are $\eta-1$, $2\eta-4$, $10\eta-22$, $14-6\eta$, $7-3\eta$ for $g=2,\dots,6$,
with sum $4\eta-6$, and dividing by $4\eta-6$ gives the frequencies.
\end{proof}

The density of this intersection is $(4\eta-6)/10=(2\sqrt5-3)/5$, in agreement with
\eqref{eq:bang} for $\alpha=(1+\sqrt5)/2$, $\beta=\sqrt5$ and $(u,v,w)=(-2,5,1)$.

\section{Two independent pairs}\label{sec:indep}

When $1,1/\alpha,1/\beta$ are linearly independent over $\Q$ the two offsets are unrelated, and
the analysis is over a rectangle.

\subsection{\texorpdfstring{$\sqrt2$}{sqrt 2} and \texorpdfstring{$\sqrt3$}{sqrt 3}}

Put $\rho=\sqrt2$, $\theta=\sqrt3$ and $S=\B(\rho)\cap\B(\theta)$. For $m\in S$ the offsets are
$a=2k-m\rho\in[0,\rho)$ and $b=3l-m\theta\in[0,\theta)$, with
$P_j=\bigcup_i[j\rho-2i,(j+1)\rho-2i)$ and $Q_j=\bigcup_i[j\theta-3i,(j+1)\theta-3i)$:
\begin{center}
\begin{tabular}{ccc}
\toprule
$j$ & $P_j\cap[0,\rho)$ & $Q_j\cap[0,\theta)$\\
\midrule
1 & $[0,\,2\rho-2)$ & $[0,\,2\theta-3)$\\
2 & $[0,\,3\rho-4)\cup[2\rho-2,\,\rho)$ & $[2\theta-3,\,\theta)$\\
3 & $[3\rho-4,\,\rho)$ & $[0,\,4\theta-6)$\\
4 & $[0,\,5\rho-6)$ & $[4\theta-6,\,\theta)$\\
\bottomrule
\end{tabular}
\end{center}
Here $3\rho-4\approx0.243$, $2\rho-2\approx0.828$ and $5\rho-6\approx1.071$. Reading off the
least $j$ with $a\in P_j$ and $b\in Q_j$ gives the gap as a function of $(a,b)$:
\begin{center}
\begin{tabular}{lccc}
\toprule
 & $b\in[0,\,2\theta-3)$ & $b\in[2\theta-3,\,4\theta-6)$ & $b\in[4\theta-6,\,\theta)$\\
\midrule
$a\in[0,\,3\rho-4)$ & 1 & 2 & 2\\
$a\in[3\rho-4,\,2\rho-2)$ & 1 & 3 & 4\\
$a\in[2\rho-2,\,\rho)$ & 3 & 2 & 2\\
\bottomrule
\end{tabular}
\end{center}
For example, in the middle cell $b\notin Q_1$, $a\notin P_2$, and $a\in P_3$, $b\in Q_3$; in
the cell to its right $b\notin Q_3$ as well, and $a\in P_4$, $b\in Q_4$.

Since $a/2$ and $b/3$ are the fractional parts of $-m/\rho$ and $-m/\theta$, and
$1,1/\sqrt2,1/\sqrt3$ are linearly independent over $\Q$, Weyl's theorem shows that for
$m\in S$ the point $(a,b)$ is equidistributed in $[0,\rho)\times[0,\theta)$. The frequency of
a gap is therefore the total area of its cells divided by $\rho\theta$. The areas are
$(2\rho-2)(2\theta-3)$, $(2\rho-2)(3-\theta)$, $2(2-\rho)(2\theta-3)$ and
$(2-\rho)(6-3\theta)$ for $g=1,2,3,4$, which gives the frequencies in
Theorem~\ref{thm:gaps}.

\subsection{The cube roots of two and four}

Put $\lambda=2^{1/3}$ and $S=\B(\lambda)\cap\B(\lambda^2)$. By Lemma~\ref{lem:hit} with
$\sigma=\lambda^2$ for $\B(\lambda)$ and $\sigma=\lambda$ for $\B(\lambda^2)$, the offsets of
$m\in S$ are $a=2k-m\lambda^2\in[0,\lambda^2)$ and $b=2l-m\lambda\in[0,\lambda)$, with
$P_j=\bigcup_i[j\lambda^2-2i,(j+1)\lambda^2-2i)$ and
$Q_j=\bigcup_i[j\lambda-2i,(j+1)\lambda-2i)$:
\begin{center}
\begin{tabular}{ccc}
\toprule
$j$ & $P_j\cap[0,\lambda^2)$ & $Q_j\cap[0,\lambda)$\\
\midrule
1 & $[0,\,2\lambda^2-2)$ & $[0,\,2\lambda-2)$\\
2 & $[0,\,3\lambda^2-4)\cup[2\lambda^2-2,\,\lambda^2)$ & $[2\lambda-2,\,\lambda)$\\
3 & $[3\lambda^2-4,\,\lambda^2)$ & $[0,\,4\lambda-4)$\\
4 & $[4\lambda^2-6,\,\lambda^2)$ & $[4\lambda-4,\,\lambda)$\\
\bottomrule
\end{tabular}
\end{center}
Here $\lambda\approx1.2599$, $\lambda^2\approx1.5874$, $3\lambda^2-4\approx0.762$,
$2\lambda^2-2\approx1.175$, $4\lambda^2-6\approx0.350$, $2\lambda-2\approx0.520$ and
$4\lambda-4\approx1.040$. The gap as a function of $(a,b)$ has the same shape as before:
\begin{center}
\begin{tabular}{lccc}
\toprule
 & $b\in[0,\,2\lambda-2)$ & $b\in[2\lambda-2,\,4\lambda-4)$ & $b\in[4\lambda-4,\,\lambda)$\\
\midrule
$a\in[0,\,3\lambda^2-4)$ & 1 & 2 & 2\\
$a\in[3\lambda^2-4,\,2\lambda^2-2)$ & 1 & 3 & 4\\
$a\in[2\lambda^2-2,\,\lambda^2)$ & 3 & 2 & 2\\
\bottomrule
\end{tabular}
\end{center}
Now $a/2$ and $b/2$ are the fractional parts of $-m/\lambda$ and $-m/\lambda^2$, and
$1,1/\lambda,1/\lambda^2$ are linearly independent over $\Q$ because $\lambda$ has degree $3$.
The areas are $(2\lambda^2-2)(2\lambda-2)$, $(2\lambda^2-2)(2-\lambda)$,
$2(2-\lambda^2)(2\lambda-2)$ and $(2-\lambda^2)(4-3\lambda)$, and dividing by
$\lambda^2\cdot\lambda=2$ gives the frequencies. This completes the proof of
Theorem~\ref{thm:gaps}.

\section{The drift of \texorpdfstring{\seqnum{A352676}}{A352676}}\label{sec:drift}

We return to $\theta=\sqrt3$ and $S=\B(\theta)\cap\B(1+\theta)$, and write
\[
\xi=\theta-1,\qquad w(m)=\fl{m\theta},\qquad N(M)=\#\{m\in S: m\le M\}.
\]
Note that $\fl{m\xi}=w(m)-m$ and $(3/2)(1+\theta)=3/\xi$. Define
\[
F(M)=3N(M)-\fl{(M+1)\xi}\quad(M\ge0),\qquad G(n)=F(a(n))=3n-\fl{(a(n)+1)\xi}.
\]
We first show that this $G$ is the walk of Theorem~\ref{thm:walk}.

\begin{lemma}[Step rule]\label{lem:step}
Let $m\in S$ and let $m+g$ be the next element of $S$. Then
$w(m+g+1)-w(m+1)$ equals $g+2$ if $g\in\{2,3\}$ and $g+4$ if $g\in\{5,6\}$. Consequently
$G(n+1)-G(n)$ is $+1$ after a short gap and $-1$ after a long gap, and $G(1)=-1$.
\end{lemma}

\begin{proof}
With $a=3k-m\theta$ as in Section~\ref{sec:sqrt3} we have $(m+j)\theta=3k+j\theta-a$, so
$w(m+j)=3k+\fl{j\theta-a}$. The values of $\fl{\theta-a}$ and $\fl{(g+1)\theta-a}$ are constant
on each cell and are listed in the last column of Table~\ref{tab:sqrt3cells}; their difference is
$4,5,9,10$ for $g=2,3,5,6$. Since $\fl{(M+1)\xi}=w(M+1)-M-1$,
\[
G(n+1)-G(n)=3-\bigl(w(a(n)+g+1)-w(a(n)+1)\bigr)+g ,
\]
which is $1$ for short gaps and $-1$ for long gaps. Finally $a(1)=5$ and
$G(1)=3-\fl{6\xi}=3-4=-1$.
\end{proof}

\begin{proof}[Proof of Theorem~\ref{thm:walk}(a)]
Let $M=a(n)$ and $\varphi_n=\{(M+1)\xi\}=\{(M+1)\theta\}$. By definition
$3n=G(n)+(M+1)\xi-\varphi_n$, so
\[
\frac{3n}{\xi}=M+1+\frac{G(n)-\varphi_n}{\xi}
\quad\text{and}\quad
d(n)=M-\Bigl\lfloor\frac{3n}{\xi}\Bigr\rfloor=-1-\Bigl\lfloor\frac{G(n)-\varphi_n}{\xi}\Bigr\rfloor .
\]
Since $M\in\B(1+\theta)$, Lemma~\ref{lem:member} gives $\fl{(M+1)\xi/2}>\fl{M\xi/2}$, hence
$\fl{(M+1)\xi}>\fl{M\xi}$, hence $\varphi_n=\{M\xi\}+\xi-1<\xi$; and $\varphi_n>0$ by
irrationality. If $G(n)=0$ then $(G(n)-\varphi_n)/\xi\in(-1,0)$ and $d(n)=0$. If $G(n)\ge1$
then $(G(n)-\varphi_n)/\xi>1/\xi-1>0$ and $d(n)\le-1$. If $G(n)\le-1$ then
$(G(n)-\varphi_n)/\xi<-1/\xi<-1$ and $d(n)\ge1$.
\end{proof}

It remains to evaluate $F$ at the numbers $q_k$. For $m\ge1$ let
\[
c(m)=w(m+1)-w(m)-1=\fl{(m+1)\xi}-\fl{m\xi}\in\{0,1\},\qquad
\kappa(m)=\bigl(w(m)+3m\bigr)\bmod 6 .
\]

\begin{lemma}\label{lem:kappa}
For $m\ge1$, $m\in S$ if and only if $c(m)=1$ and $\kappa(m)\in\{1,5\}$. Consequently
\[
F(M)=\sum_{\substack{1\le m\le M\\ c(m)=1}} f(\kappa(m)),\qquad
f(\kappa)=\begin{cases}2,&\kappa\in\{1,5\},\\-1,&\kappa\in\{0,2,3,4\},\end{cases}
\]
and $\kappa(1)=4$, $\kappa(m+1)\equiv\kappa(m)+4+c(m)\pmod 6$.
\end{lemma}

\begin{proof}
Since $m/\theta=m\theta/3$ and $\fl{x/3}=\fl{\fl{x}/3}$, we have $\fl{m/\theta}=\fl{w(m)/3}$,
so by Lemma~\ref{lem:member}, $m\in\B(\theta)$ if and only if $\fl{w(m+1)/3}>\fl{w(m)/3}$.
Similarly $m/(1+\theta)=m\xi/2$ and $\fl{m\xi/2}=\fl{(w(m)-m)/2}$, so $m\in\B(1+\theta)$ if and
only if $\fl{(w(m+1)-m-1)/2}>\fl{(w(m)-m)/2}$. If $c(m)=0$ then $w(m+1)-m-1=w(m)-m$ and
$m\notin\B(1+\theta)$. If $c(m)=1$ then $w(m+1)=w(m)+2$, so $m\in\B(1+\theta)$ if and only if
$w(m)-m$ is odd, and $m\in\B(\theta)$ if and only if $3\nmid w(m)$. These two conditions say
that $w(m)+3m$ is odd and not divisible by $3$, that is, $\kappa(m)\in\{1,5\}$. The formula for
$F$ follows because $\sum_{m\le M}c(m)=\fl{(M+1)\xi}$. The recurrence for $\kappa$ is immediate
from $w(m+1)=w(m)+1+c(m)$, and $\kappa(1)=1+3=4$.
\end{proof}

The word $c(1)c(2)c(3)\cdots$ is the characteristic Sturmian word of slope
$\xi=[0;1,2,1,2,\dots]$. Define words $s_k$ by
\[
s_{-1}=1,\quad s_0=0,\quad s_1=1,\qquad
s_k=\begin{cases}s_{k-1}s_{k-1}s_{k-2},&k\ge2\text{ even},\\ s_{k-1}s_{k-2},&k\ge3\text{ odd},\end{cases}
\]
so that $s_2=110$, $s_3=1101$, $s_4=11011101110$, and $s_k$ has length $q_k$. By the theory of
standard words \cite[Section~2.2.2]{lothaire}, for every $k\ge1$ the word $s_k$ is the prefix
of
$c(1)c(2)\cdots$ of length $q_k$.

For a word $v=v_1\cdots v_\ell$ over $\{0,1\}$ let $\gamma(v)\in\Z/6$ be the class of $4\ell$
plus the number of letters $1$ in $v$, and for $\kappa\in\Z/6$ let
\[
\Phi_v(\kappa)=\sum_{i:\,v_i=1} f\bigl(\kappa+\gamma(v_1\cdots v_{i-1})\bigr).
\]
Then $\Phi_{vv'}(\kappa)=\Phi_v(\kappa)+\Phi_{v'}(\kappa+\gamma(v))$ and
$\gamma(vv')=\gamma(v)+\gamma(v')$, and Lemma~\ref{lem:kappa} gives
$F(M)=\Phi_{c(1)\cdots c(M)}(4)$. We write $\Phi_k=\Phi_{s_k}$ as the vector
$(\Phi_k(0),\dots,\Phi_k(5))$ and $\gamma_k=\gamma(s_k)$.

\begin{proposition}\label{prop:blocks}
For $k=4j+r\ge0$ with $0\le r\le3$,
\begin{center}
\begin{tabular}{clc}
\toprule
$r$ & $\Phi_k$ & $\gamma_k$\\
\midrule
$0$ & $j\,(1,-2,-2,-2,1,4)$ & $4$\\
$1$ & $(-1,2,-1,-1,-1,2)+j\,(-1,-1,-1,-1,2,2)$ & $5$\\
$2$ & $(1,1,1,-2,-2,1)+j\,(2,2,-1,-4,-1,2)$ & $2$\\
$3$ & $(0,0,0,0,-3,3)+j\,(1,1,1,-2,-2,1)$ & $1$\\
\bottomrule
\end{tabular}
\end{center}
In particular $F(q_{4j+1})=\Phi_{4j+1}(4)=2j-1$ and $F(q_{4j+3})=\Phi_{4j+3}(4)=-2j-3$.
\end{proposition}

\begin{proof}
The cases $k=0,1$ are direct: $s_0=0$ gives $\Phi_0=0$, $\gamma_0=4$, and $s_1=1$ gives
$\Phi_1=(f(0),\dots,f(5))$, $\gamma_1=5$. For $k\ge2$ the definition of $s_k$ gives
\[
\Phi_k(\kappa)=\begin{cases}
\Phi_{k-1}(\kappa)+\Phi_{k-1}(\kappa+\gamma_{k-1})+\Phi_{k-2}(\kappa+2\gamma_{k-1}),&k\text{ even},\\
\Phi_{k-1}(\kappa)+\Phi_{k-2}(\kappa+\gamma_{k-1}),&k\text{ odd},
\end{cases}
\]
and $\gamma_k=2\gamma_{k-1}+\gamma_{k-2}$ or $\gamma_{k-1}+\gamma_{k-2}$ respectively, so
$\gamma_k$ runs through $4,5,2,1$ periodically. The claim follows by induction: both sides of
the recurrence are affine in $j$, and one compares constant terms and coefficients of $j$ in the
four cases $k\equiv2,3,0,1\pmod 4$. For $k=4j+2$, for instance, $\gamma_{k-1}=5$, and writing
$\tau^e u$ for the vector with entries $u(\kappa+e)$,
\begin{align*}
(-1,2,-1,-1,-1,2)+\tau^5(-1,2,-1,-1,-1,2)&=(1,1,1,-2,-2,1),\\
(-1,-1,-1,-1,2,2)+\tau^5(-1,-1,-1,-1,2,2)\qquad&\\
{}+\tau^{4}(1,-2,-2,-2,1,4)&=(2,2,-1,-4,-1,2),
\end{align*}
as required. The other three cases are identical in form; all four are checked by the program
\texttt{verify\_induction.py} in the repository.
\end{proof}

\begin{proof}[Proof of Theorem~\ref{thm:walk}(b) and of Theorem~\ref{thm:a352676}(iii)]
Let $a(n)\le M<a(n+1)$. Then $G(n)-F(M)=\fl{(M+1)\xi}-\fl{(a(n)+1)\xi}$ is the number of
$m$ with $a(n)<m\le M$ and $c(m)=1$. With $g=a(n+1)-a(n)$, the number of $m$ with
$a(n)<m\le a(n+1)$ and $c(m)=1$ is $w(a(n+1)+1)-w(a(n)+1)-g$, which by
Lemma~\ref{lem:step} is $2$ for a short gap and $4$ for a long gap; one of these is
$m=a(n+1)$ itself. Hence $F(M)\le G(n)\le F(M)+3$. Taking $M=q_{4j+1}$ and
$M=q_{4j+3}$ and using Proposition~\ref{prop:blocks} gives (b).

So $G$ is unbounded above and below. As $G(n+1)-G(n)=\pm1$, it changes sign infinitely often and
takes the value $0$ each time. By Theorem~\ref{thm:walk}(a), $d(n)=0$ infinitely often, and
since $|d(n)|\ge|G(n)|/\xi-1$ with sign opposite to that of $G(n)$, $d$ is unbounded above and
below.
\end{proof}

\begin{remark}
Begunts and Goryashin \cite{bg25} prove that the counting function of this intersection differs
from its main term by $O(\log^5 N)$. Theorem~\ref{thm:walk}(b) shows that the difference is not
bounded, and is at least of order $\log N$ infinitely often in each direction.
\end{remark}

\section{Formal verification}\label{sec:lean}

The following statements are proved in Lean~4 with Mathlib. The development defines
\[
\texttt{Beatty}\ r=\{m\in\Z:\ \exists k\in\Z,\ k>0\text{ and }\fl{kr}=m\}
\]
and says that $g$ is a gap of a set $S\subseteq\Z$ (\texttt{IsGap S g}) when there is $m\in S$
with $m+g\in S$ and $m+j\notin S$ for all $0<j<g$.

\begin{itemize}
\item For each of the four pairs, a theorem \texttt{gaps} stating that after any element $m$ of
the intersection the next element is $m+g$ for some $g$ in the gap set of
Theorem~\ref{thm:gaps}, and a theorem \texttt{isGap\_iff} stating that for $g>0$,
\texttt{IsGap S g} holds if and only if $g$ is in that set. The names are
\texttt{BeattyGaps.Sqrt3.isGap\_iff}, \texttt{BeattyGaps.Golden.isGap\_iff},
\texttt{BeattyGaps.Sqrt2Sqrt3.isGap\_iff} and \texttt{BeattyGaps.CubeRoot2.isGap\_iff}. The
last is proved for every real $c>0$ with $c^3=2$.
\item \texttt{BeattyGaps.Sqrt3.four\_not\_gap}: $4$ is not a gap of \seqnum{A352676}.
\item \texttt{BeattyGaps.Sqrt3.walk\_step} and \texttt{next\_elem}: Lemma~\ref{lem:step}, in
the form that $F$ changes by exactly $+1$ across a short gap and $-1$ across a long gap.
\item \texttt{inS\_iff} and \texttt{F\_eq\_Phi}: Lemma~\ref{lem:kappa} and the identity
$F(M)=\Phi_{c(1)\cdots c(M)}(4)$.
\item \texttt{good\_all}: Proposition~\ref{prop:blocks}.
\item \texttt{drift\_of\_F} and \texttt{drift\_eq\_zero\_iff}: at an element of the intersection,
$d=0$ if and only if $F=0$; if $F\ge1$ then $d\le-F$; if $F\le-1$ then $d\ge-F$. This is the part
of Theorem~\ref{thm:walk}(a) needed below; the exact floor formula is not formalized.
\item The theorems \texttt{d\_zero\_infinitely\_often}, \texttt{d\_unbounded\_above}
and\\ \texttt{d\_unbounded\_below}: Theorem~\ref{thm:a352676}(iii), for $a(n)$ defined as the $n$-th
element and $d(n)=a(n)-\fl{(3/2)(1+\sqrt3)\,n}$, \emph{assuming} the hypothesis
\texttt{BlockHypothesis}.
\end{itemize}

The hypothesis \texttt{BlockHypothesis} is the statement, quoted in Section~\ref{sec:drift} from
the theory of standard words, that for every $k\ge1$ the first $q_k$ letters of
$c(1)c(2)\cdots$ form the word $s_k$. It is not proved in Lean; Mathlib has no theory of
Sturmian words. Lean does verify it for $k\le4$ (\texttt{block\_small}), which guards against an
error in the formal definitions.

The proofs of the gap sets follow Sections~\ref{sec:prelim}--\ref{sec:indep}. A program computes
the cells and, for each cell, a certificate: the gap $g$, the integers witnessing that $m+g$ lies
in both sequences, and for each $j<g$ the sequence that misses $m+j$ together with the two
consecutive multiples between which the relevant interval falls. Lean checks each certificate by
linear arithmetic over the reals, treating the irrational number as an unknown subject to a
numerical enclosure such as $1.732<\sqrt3<1.7321$. An error in the program can therefore cause a
proof to fail but cannot cause a false statement to be accepted. The formalization of
Section~\ref{sec:drift} is written by hand and follows the text. All proofs use no axioms beyond
the three standard ones of Lean's classical logic.

Not formalized: the frequencies in Theorem~\ref{thm:gaps} and the statement that each gap occurs
infinitely often (Lean shows that each occurs at least once); Theorem~\ref{thm:a352676}(ii); the
hypothesis \texttt{BlockHypothesis}; and the precise constants in Theorem~\ref{thm:walk}. As an
independent check, $d(n)$ was computed exactly for all elements up to $2\cdot10^9$, where it
ranges from $-22$ to $21$, and the block sums of Proposition~\ref{prop:blocks} agree with
brute-force computation for all $q_k\le 600000$.

\section{Open questions}\label{sec:open}

\begin{enumerate}
\item The conjectures of type (3) in \seqnum{A346308} and \seqnum{A347793}, and the related
conjecture in \seqnum{A347469}, remain open. For these pairs the deviation is a sum over a
two-dimensional torus translation and we know of no analogue of the step rule.
\item Is $d(n)=O(\log n)$ for \seqnum{A352676}?
\item The method of Sections~\ref{sec:sqrt3}--\ref{sec:indep} determines the gap set of
$\B(\alpha)\cap\B(\beta)$ whenever the comparisons between the endpoints can be decided, and
the finiteness theorem of Fraenkel and Holzman guarantees that it terminates. Is there an
explicit bound for the largest gap in terms of $\alpha$ and $\beta$?
\item The two independent pairs treated here have the same table of gaps. For which pairs does
this happen?
\end{enumerate}

\section{Use of AI tools}\label{sec:ai}

This work was carried out with an AI system, Claude (Anthropic), under the direction of the
author. The AI system proposed the topic after a literature search, found and carried out the
proofs, wrote the Lean formalization and all computer programs, and wrote the text of this
draft. The author directed the work, chose among the options presented, and obtained sources.
The statements listed in Section~\ref{sec:lean} as formalized have been checked by Lean; the
remaining arguments have been checked by exact computation as described there but have not yet
been reviewed by a mathematician. The literature search is recorded in the repository.

\begin{thebibliography}{99}

\bibitem{bang}
T. Bang, On the sequence $[n\alpha]$, $n=1,2,\dots$. Supplementary note to the preceding paper
by Th.~Skolem, \emph{Math. Scand.} \textbf{5} (1957), 69--76.

\bibitem{bg22}
A. V. Begunts and D. V. Goryashin, On the intersection of two homogeneous Beatty sequences,
\emph{Chebyshevskii Sb.} \textbf{23} (2022), 145--151.

\bibitem{bg25}
A. V. Begunts and D. V. Goryashin, On the intersection of two dependent Beatty sequences,
\emph{Mat. Zametki} \textbf{118} (2025), 670--678; English translation in
\emph{Math. Notes} \textbf{118} (2025), 933--939.

\bibitem{fh}
A. S. Fraenkel and R. Holzman, Gap problems for integer part and fractional part sequences,
\emph{J. Number Theory} \textbf{50} (1995), 66--86.

\bibitem{kalmynin}
A. Kalmynin, Answer to ``Limit associated with two Beatty sequences that are not a Beatty
pair'', MathOverflow, 2022, \url{https://mathoverflow.net/questions/418749}.

\bibitem{kn}
L. Kuipers and H. Niederreiter, \emph{Uniform Distribution of Sequences}, Wiley, 1974.

\bibitem{lothaire}
M. Lothaire, \emph{Algebraic Combinatorics on Words}, Cambridge University Press, 2002.

\bibitem{mathlib}
The mathlib Community, The Lean mathematical library, in \emph{Proceedings of the 9th ACM
SIGPLAN International Conference on Certified Programs and Proofs}, ACM, 2020, pp.~367--381.

\bibitem{lean4}
L. de Moura and S. Ullrich, The Lean 4 theorem prover and programming language, in
\emph{Automated Deduction -- CADE 28}, Lecture Notes in Comput. Sci., Vol.~12699, Springer,
2021, pp.~625--635.

\bibitem{simpson}
J. Simpson, Intersecting rational Beatty sequences, \emph{Integers} \textbf{13} (2013),
Paper A50.

\bibitem{skolem}
Th. Skolem, On certain distributions of integers in pairs with given differences,
\emph{Math. Scand.} \textbf{5} (1957), 57--68.

\bibitem{oeis}
N. J. A. Sloane et al., The On-Line Encyclopedia of Integer Sequences, 2026,
\url{https://oeis.org}.

\bibitem{weyl}
H. Weyl, \"Uber die Gleichverteilung von Zahlen mod.\ Eins, \emph{Math. Ann.} \textbf{77}
(1916), 313--352.

\end{thebibliography}

\bigskip
\noindent 2020 \emph{Mathematics Subject Classification}: Primary 11B83; Secondary 11J71, 68R15,
68V20.

\noindent \emph{Keywords}: Beatty sequence, gap, Sturmian word, discrepancy, formal proof.

\begin{sloppypar}
\noindent (Concerned with sequences
\seqnum{A000201}, \seqnum{A001951}, \seqnum{A022838}, \seqnum{A022839}, \seqnum{A054088},
\seqnum{A346308}, \seqnum{A347469}, \seqnum{A347793}, \seqnum{A351415}, \seqnum{A352674} and
\seqnum{A352676}.)
\end{sloppypar}

\end{document}
